\begin{frame}[fragile]{실습: CartPole에서 무작위 정책 실행}
\begin{lstlisting}
import gymnasium as gym

env = gym.make("CartPole-v1")  # 막대를 쓰러뜨리지 않고 균형 잡기
obs, info = env.reset(seed=0)
print("초기 상태(카트 위치, 속도, 막대 각도, 각속도):", obs)

for _ in range(3):
    action = env.action_space.sample()  # 지금은 무작위 정책
    obs, reward, terminated, truncated, info = env.step(action)
    print("행동:", action, "-> 다음 상태:", obs, "보상:", reward)

env.close()
\end{lstlisting}
  {\footnotesize 지금은 \texttt{env.action\_space.sample()}로
  \textbf{완전 무작위}로 행동 --- 2$\sim$11주차을 거치며 이 한 줄이
  ``알고리즘이 고른 행동''으로 \textbf{점점 똑똑해}지는 것을 확인.
  이 CartPole은 표 기반 알고리즘을 \textbf{시각적으로} 이해하는 데,
  그리고 9주차의 DQN에서 \textbf{다시} 등장.}
\end{frame}
